% Experimental research record; does not claim successful novelty.
\section{物理引导条件动作生成：对比与消融诊断记录}
\label{sec:physics-guided-diagnostic-20261011}
\subsection{目的、机制与证据边界}
本轮检验物理处理是否改善参考动作的可用性，同时保留原始条件。冻结DiP生成权重和CLoSD人体物理跟踪器，使用有限次条件纠偏构造候选：根据执行终点相对于固定目标的位置与朝向误差，修正下一轮生成条件，最多三个候选。每次位置调整不超过5厘米、累计不超过15厘米，朝向调整不超过5度；文本、动作长度和目标时刻不变。
物理精化、执行状态反馈和闭环数据扩展已有直接先例\cite{closd,parc,claims}，不能将这些机制本身作为创新。本轮是CLoSD发布组件的适配诊断，不是Morph、PhysDiff或CLAIMS完整复现。
\subsection{数据、对照与验收}
跳跃与舞蹈各选64个训练来源、12个验证来源、24个测试来源。按官方原始MoCap来源分组，镜像与切片不跨本轮划分；发布权重的预训练暴露并未排除。测试48个来源、三个目标水平偏移和三个生成种子，每方法432次输入。任务仅由文本标签预先筛选，4秒窗口不保证包含完整技能事件，因此不能称为完整舞蹈或跳跃专项基准。
训练来源预留给后续数据效用验证；当前未重训生成器或跟踪器。
后续独立复核发现，参考及目标锚点来自角色创建位置，而实际角色在开始控制前经过重置，二者没有同步平移。同一案例和种子内各方法共用原目标，但种子变化同时改变任务几何，三个目标偏移案例也受不同锚点影响。因此，本轮不能视为受控的15厘米偏移实验，不能据三个种子推定随机种子稳健性；计划前缀与执行前缀的比较亦可能受初始化错位影响。这是主要测量有效性缺陷，而非可以忽略的坐标表示差异。
对照包括未经精化的生成参考、普通固定跟踪精化、等预算随机筛选及执行前缀反馈。消融分别去除条件误差反馈、跨原始来源打乱反馈、仅按跟踪表现选候选，以及同时改用计划前缀和计划坐标锚点。所有参考经同一冻结跟踪器独立回放；构建失败、提前终止及不合格候选保留。
联合验收同时要求条件、物理代理、跟踪与完整执行合格。足部滑移按支撑状态、足部碰撞表面及角速度计算，是代理而非真实接触点测量；腾空本身不判为漂浮。验证阶段用于检查测量有效性，数值阈值已在验证前写入，未按验证分布重新校准。这与原计划的阈值校准程序存在偏差；本轮不在看到测试结果后放宽阈值。
\subsection{独立回放结果}
\begin{center}\small
\begin{tabular}{p{6.5cm}rrr}\hline
方法 & 尝试 & 联合合格 & 完成回放 \\\hline
未精化生成参考 & 432 & 1 & 279\\
普通物理精化 & 432 & 0 & 263\\
等预算随机候选筛选 & 432 & 0 & 265\\
原生执行前缀反馈（单候选） & 432 & 0 & 280\\
原生执行前缀反馈（等预算筛选） & 432 & 0 & 279\\
候选方法：执行前缀＋有界条件反馈 & 432 & 0 & 276\\
消融：去除条件误差反馈 & 432 & 0 & 280\\
消融：打乱原始来源反馈 & 432 & 0 & 280\\
消融：计划前缀＋计划锚点 & 432 & 0 & 260\\
消融：仅按跟踪表现筛选 & 432 & 0 & 280\\
\hline\end{tabular}\end{center}

回放完成表示达到120帧且未触发提前终止，不等同于条件与物理联合合格；提前终止也不能全部称为跌倒。
\subsection{配对比较与消融结论}
普通物理精化：候选方法提升 0.00 个百分点，来源聚类95\%置信区间[0.00, 0.00]。
等预算随机候选筛选：候选方法提升 0.00 个百分点，来源聚类95\%置信区间[0.00, 0.00]。
原生执行前缀反馈（等预算筛选）：候选方法提升 0.00 个百分点，来源聚类95\%置信区间[0.00, 0.00]。
消融：去除条件误差反馈：候选方法提升 0.00 个百分点，来源聚类95\%置信区间[0.00, 0.00]。
消融：打乱原始来源反馈：候选方法提升 0.00 个百分点，来源聚类95\%置信区间[0.00, 0.00]。
消融：计划前缀＋计划锚点：候选方法提升 0.00 个百分点，来源聚类95\%置信区间[0.00, 0.00]。
消融：仅按跟踪表现筛选：候选方法提升 0.00 个百分点，来源聚类95\%置信区间[0.00, 0.00]。

进入后续阶段要求候选方法对每个等预算对照均提高至少10个百分点、来源聚类置信区间排除零、三个种子同向，并通过条件、语义与多样性保护。
\subsection{阶段判定与论文贡献分类}
本轮人体门槛未通过，因此按停止规则不开展G1与下游九组训练，不补造未执行结果。结论分为：
\begin{itemize}
\item 有效且有技术差异：本轮未建立此类结论。
\item 有效但属于已有机制：物理精化与执行反馈是已有路线；本轮不能由回放完成率将其直接解释为联合质量收益。
\item 未获支持：本轮有界条件纠偏未达到预定联合验收提升门槛，不能作为已验证的论文创新。
\end{itemize}
验证集已出现零联合合格率，存在判据交集的地板效应风险。加上初始化坐标校验失败，本轮数值不能作为候选机制无效的证据。后续必须先修复初始化与测量，再独立验证并冻结协议，不能通过放宽测试判据取得正结果。
本轮在实现复核中发生过修正，并逐条核验复用的早期输出。因此归类为经过修订的诊断实验，不能称为无修订的确认性实验。原执行请求、轨迹散列、版本快照、完整尝试与失败记录均保留，缺失的历史执行绑定不追溯补造。语义与多样性采用冻结HumanML3D文本--动作嵌入评价器，是学习型代理而非人工真值。

语义与多样性代理的实际状态：pass。

语义与多样性代理基于独立回放的实际状态，不能代替生成参考本身的独立语义验收。

正式记录含74个批次、10512次构建或回放首轮尝试，其中4320次独立回放。保守计费GPU时间为 1.959 小时，低于72小时上限；本轮输出限制15GiB，旧实验与其他项目未删除。
完整数值、对照、消融、版本与失败记录见本地实验报告和\href{https://github.com/Pendragon08/conditional-human-motion-generation}{现有GitHub项目}。


% RESET_REPAIR_RECORD_20261011
\subsection{独立初始化修复与停止依据}
独立初始化复核计划216例，实际启动12例、完成12例，测量校验未通过。首动作前实际与参考根节点的最大XY分量误差为 0.00071254373 米，引导重置与首动作前的刚体位置最大分量差为 0.03500697 米；其余导出状态逐项差异见摘要，不能将不同单位的差异合为一个误差指标。已运行来源的三个目标偏移最大误差为 3.0636787e-06 米。跨种子、跨模式和完整验证未运行。外层进程在下一批前触发保护其他工作负载的资源断言，退出码为1；原始状态文件残留RUNNING，不代表仍在运行。具体失败项见独立测量摘要，不据此开展机制有效性比较。本复核不改变质量阈值，不与旧轮合并，不包含候选反馈的完整对比或下游训练；G1及九组策略训练按停止条件未启动。
\begin{center}\small\begin{tabular}{lrrrrrr}\hline 模式 & 尝试 & 完成 & 条件 & 物理代理 & 跟踪 & 联合 \\ \hline
计划前缀（部分测量） & 12 & 12 & 4 & 0 & 12 & 0 \\
\hline\end{tabular}\end{center}
前两次修复在控制前分别因目标未初始化及坐标断言失败，原失败保留。V3的XY锚点改善不等于完整初态通过；首批导出刚体状态仍有差异，后续批次未启动。所有质量阈值保持不变。保守合计 1.989 GPU小时。
